\subsection{Nuclear LWR, rest of world (11,000~\$/kW)}\label{sec:nuclear-lwr-row}

This bin covers new water-cooled reactors built outside China and India.
Such projects have become more expensive over the past half-century \citep{eashgates2020}.

The CAPEX of 11,000~\$/kW is the average of Western nuclear projects of the last decades (Figure~\ref{fig:nuclear-strip}).
It is slightly below the 11,900~\$/kW of technology-data \citep{technologydata}, the usual source for European studies.
The most recent Western reactors, Vogtle, Flamanville and Hinkley Point~C, came out higher still at 12,000--22,000~\$/kW.
Either way, this cost is too high for the model to pick up water-cooled nuclear in the West, even in a carbon-neutral setting.

Some countries have built nuclear plants at reliably lower cost than the West, notably South Korea \citep{lovering2016} and Russia.
We neglect this nuance because neither lies within our regions of focus.

The learning rate is zero.
The historical record for large light-water reactors ranges from negative to 6\,\% \citep{rubin2015,eashgates2020}.

All other parameters are those of the China--India bin; the lead time is 10~years.
The build limit of 7~GW/yr is the most construction started outside China and India in any year of the last decade (2024) \citep{gem2026}; at this CAPEX, the model is unlikely to reach it.
